\uuid{0zIe}
\exo7id{857}
\titre{exo7 857}
\auteur{gourio}
\organisation{exo7}
\datecreate{2001-09-01}
\isIndication{false}
\isCorrection{false}
\chapitre{Equation différentielle}
\sousChapitre{Résolution d'équation différentielle du premier ordre}
\cpge{mpsi-mp2i.fonctions-valeurs-reelles-complexes.d, pcsi.techniques-calcul-fonctions.d, ptsi.techniques-calcul-fonctions.d, mpsi-mp2i.limites-continuite.a, pcsi.limites-continuite.a, ptsi.limites-continuite.a, bcpst1.limites-continuite.a, bcpst1.derivation.a}
\module{Analyse}
\niveau{L2}
\difficulte{}

\contenu{
\question{
Soit $f\in C^{1}({\Rr},\Cc),\alpha \in {\Rr}^{+*}. $ Montrer que
si :
$$\lim\limits_{x\rightarrow \infty }\left( f^{\prime }(x)+\alpha f(x)\right)=0 $$
alors :
$$\lim_{x\rightarrow \infty }f^{\prime }(x)=\lim_{x\rightarrow \infty }f(x)=0. $$
}
}
